% Definition file for row def_3_1 -- "CDMG".
% Auto-generated stub by scaffold/solve_chapter.py at row start. The
% manager agent (or a worker dispatched by it) fills the body below with the
% Lean-faithful definition statement(s). Stays close to the lecture notes.
%
% This file is a sibling of `main.tex` inside `<subsection>/tex/`. The
% `subfiles` package lets it compile both standalone (resolving its
% preamble via `main.tex`) and as part of `main.tex` via `\subfile{...}`.

\documentclass[main]{subfiles}

\begin{document}
% ref: def_3_1
% title: CDMG
\def\rowref{def\_3\_1}
\phantomsection\label{def_3_1}

\begin{Def}[Conditional directed mixed graphs (CDMG)]
    \label{def-cdmg}
    A \emph{conditional directed mixed graph (CDMG)} $G$---per definition---consists of two finite, disjoint sets of
    vertices (also called nodes):
    \begin{enumerate}[label=\roman*.)]
        \item $J$, whose elements are called input nodes,
        \item $V$, whose elements are called output nodes,
    \end{enumerate}
    with both $J$ and $V$ finite and $J \cap V = \emptyset$;
%    such that $J \cup V \neq \emptyset$;
    and two sets of edges (here ``disjoint'' is to be read as trivial type-level disjointness only: directed edges and bidirected edges are objects of distinct types, and no graph-theoretic mutual-exclusion is imposed between $E$ and $L$):
    \begin{enumerate}[resume,label=\roman*.)]
        \item $E \subseteq (J \cup V) \times V$, the set of directed edges,
        \item $L \subseteq V \times V$, the set of bidirected edges, subject to the irreflexivity and symmetry constraint
            \[ \forall\, (v_1, v_2) \in V \times V :\quad (v_1, v_2) \in L \;\Longrightarrow\; v_1 \neq v_2 \,\land\, (v_2, v_1) \in L. \]
            Equivalently, $L$ may be identified with a subset of the quotient $(V \times V)/{\sim}$ under the swap relation $(v_1, v_2) \sim (v_2, v_1)$ (i.e.\ a set of unordered pairs of distinct elements of $V$); under the ordered-pair encoding above, membership of $(v_1, v_2)$ and of $(v_2, v_1)$ in $L$ always coincide, and the irreflexivity condition $v_1 \neq v_2$ applies under either encoding.
        \item[] In particular, for $v_1, v_2 \in V$ with $v_1 \neq v_2$, it is permitted to have $(v_1, v_2) \in E$ and the unordered pair $\{v_1, v_2\}$ in $L$ simultaneously; under the ordered-pair encoding of $L$, the same ordered pair $(v, w) \in V \times V$ may belong to both $E$ and $L$.
    \end{enumerate}
\end{Def}

\end{document}
